%----------------------------------------------------------------------------
% Course Syllabus for Chemistry 350 (c) 2016 by Dale J. Brugh

% Course Syllabus for Chemistry 350 is released under a
% Creative Commons Attribution-ShareAlike 4.0 International License.

% See http://creativecommons.org/licenses/by-sa/4.0/ for
% a description of your rights under this license. 
%----------------------------------------------------------------------------

\documentclass[letterpaper,oneside,onecolumn,11pt,article]{memoir}
\usepackage[T1]{fontenc}            % use T1 font encoding
\usepackage{textcomp}
\usepackage{courier}                % set courier as typewriter font
\usepackage{times}                  % set times as text font
\usepackage[scaled=0.92]{helvet}    % set Helvetica as the sans-serif font
\usepackage{mtpro2}

\usepackage{setspace}
\usepackage{amsmath}
\usepackage{graphicx,color}
\usepackage{wallpaper}
\usepackage{textcomp}
\usepackage{relsize,fancyvrb}
\usepackage{verbatim}
\usepackage{caption}
\usepackage{paralist}

\usepackage{boxedminipage}

\usepackage[bookmarks=true]{hyperref}
\hypersetup{
    unicode=false,          % non-Latin characters in Acrobat’s bookmarks
    pdftoolbar=true,        % show Acrobat’s toolbar?
    pdfmenubar=true,        % show Acrobat’s menu?
    pdffitwindow=true,      % page fit to window when opened
    pdftitle={Course Syllabus for Chemistry 350}, 
    pdfauthor={Dale J. Brugh},     % author
    pdfsubject={Physical Chemistry},   % subject of the document
    pdfnewwindow=true,      % links in new window
    pdfkeywords={classes, ch350f16}, % list of keywords
    colorlinks=true,       % false: boxed links; true: colored links
    linkcolor=black,          % color of internal links
    citecolor=green,        % color of links to bibliography
    filecolor=magenta,      % color of file links
    urlcolor=black           % color of external links
}

\definecolor{nicered}{rgb}{.647,.129,.149}
\definecolor{mutedgrey}{rgb}{0.4,0.4,0.4}
\definecolor{shadecolor}{cmyk}{0,0,0.25,0.07}
\definecolor{MyDarkBlue}{rgb}{0,0.08,0.45}
\definecolor{MarginRed}{rgb}{0.8,0.0,0.0}
\definecolor{MarginBlue}{rgb}{0.2,0.0,1.0}
\definecolor{MarginGrey}{rgb}{0.4,0.4,0.4}

%\renewcommand{\chapnumfont}{\bfseries\Huge\sffamily}
%\renewcommand{\chaptitlefont}{\bfseries\Large\sffamily}

\setsecheadstyle{\bfseries\Large\sffamily\raggedright}
\setsubsecheadstyle{\bfseries\large\sffamily\raggedright}
\setsubsubsecheadstyle{\bfseries\normalsize\sffamily\raggedright}
\renewcommand \thesection{\bfseries\arabic{section}}

% \makeatletter 
% \newcommand\addRevisionData{%
% \begin{picture}(0,0)% 
%     \put(-110,-5){%
%         \tiny% 
%         {%
%         {Published \today \enspace \copyright~Dale J. Brugh 
%         }
% }% 
% }%
% \end{picture}%
% }

\flushbottom
\setstocksize{11in}{8.5in}
%\setlength{\parskip}{5pt}
\settrims{0pt}{0pt}
%\settrimmedsize{11in}{210mm}{*}
%\setlength{\trimtop}{0pt}
%\setlength{\trimedge}{\stockwidth}
%\addtolength{\trimedge}{-\paperwidth}
\settypeblocksize{8.5in}{5.0in}{*}
\setulmargins{1.25in}{*}{*}
\setlrmargins{1.25in}{*}{*}
\setmarginnotes{5mm}{4.0cm}{\onelineskip}
\setheadfoot{\onelineskip}{4\onelineskip}
\setheaderspaces{*}{\onelineskip}{*}
\checkandfixthelayout
%\setlength \fboxsep{0.1in}
\setlength \headwidth{\textwidth+\marginparwidth+\marginparsep}
%
\makepagestyle{courseinformation}
\makerunningwidth{courseinformation}{\headwidth}

\makeheadrule{courseinformation}{\headwidth}{\normalrulethickness}
\makeheadposition{courseinformation}{flushright}{flushleft}{flushleft}{flushleft}
\makeoddhead{courseinformation}%
    {\sffamily Course Syllabus: Chemistry 350 / Fall 2016}{}{\sffamily\thepage}

    %\makeevenfoot{courseinformation}{}{}{\addRevisionData}
    %\makeoddfoot{courseinformation}{}{}{\addRevisionData}
\makepagestyle{courseinformationtitle}
\makerunningwidth{courseinformationtitle}{\headwidth}
\makeheadposition{courseinformationtitle}{flushright}{flushleft}{flushleft}{flushleft}
    %\makeevenfoot{courseinformationtitle}{}{}{\addRevisionData}
    %\makeoddfoot{courseinformationtitle}{}{}{\addRevisionData}
\pagestyle{courseinformation}
\captionsetup{labelsep=colon,aboveskip=0.25cm,justification=RaggedRight,singlelinecheck=false,labelfont={bf,sf}}
%
% --- MARGIN FIGURE COMMAND ---
%
\newcommand{\marginfigures}[4]{
\marginpar{\centering
\includegraphics[width=#1]{#2}
\captionsetup{labelsep=newline,aboveskip=-0.5cm,justification=RaggedRight,singlelinecheck=false,labelfont={bf,sf}}
\captionsetup[figure]{position=bottom}
\captionof{figure}{#3}
\label{#4}}%
}%
%
% --- MARGIN NOTE COMMAND ---
%
\newcommand{\marginnote}[2]
{%
\marginpar{\raggedright\vspace{#1}\begin{Spacing}{0.65}\sffamily{{\tiny$\blacktriangleright$~\scriptsize#2}}\end{Spacing}} %
}
%
% --- SET UP TITLE ---
%
\setlength{\droptitle}{0.0in}
\backmatter
\pretitle{\noindent\huge\sffamily Course Syllabus \LARGE\par\noindent} 
\posttitle{\par\vskip 2.0em}
\preauthor{}
\postauthor{\par}
\predate{}
\postdate{\noindent\rule{\linewidth}{0.3pt}}
\title{Chemistry 350 / Fall 2016}
\date{}
\author{}
%
% --- BEGIN DOCUMENT ---
%
\begin{document}
\setsecnumdepth{subsubsection}
\maketitle
%\setsecnumdepth{subsection}
\thispagestyle{courseinformationtitle}
%
% --- INSTRUCTOR ---
%
\section{Instructor}
\begin{tabular}{rl|rl}
Name: & Dr. Dale J. Brugh & Office Phone: & 740-368-3530 \\
Email: & \href{mailto:djbrugh@owu.edu}{djbrugh@owu.edu} & Cell Phone: & 614-746-2397 \\
Office: & SCSC 262 & &  \\
\end{tabular}
%
% --- MEETINGS ---
%
\section{Meetings}
%\renewcommand{\arraystretch}{1.2}
\begin{tabular}{crcrl}
TR & 8:00 a.m. & to & 9:50 a.m. & SCSC 167\\
\end{tabular} \\[0.09in]
Review meetings may occasionally offered on Wednesday evenings from 7:00~p.m.\ to 8:30~p.m. Attendance at review meetings is optional. 
%
% --- WEBSITE ---
%
\section{Website}
The course website can be reached by logging in to OWU's Blackboard site at \href{https://bb.owu.edu}{bb.owu.edu} and navigating to the course named \href{https://bb.owu.edu/webapps/blackboard/execute/launcher?type=Course&id=_17564_1&url=}{2016FALLCHEM3501}. The website is an important extension of the syllabus and should be read carefully. 
%
% --- PREREQUISITES ---
%
\section{Prerequisites}
To \marginnote{-0.1in}{If you have not met these prerequisites, see me immediately.}take this course you must have passed two semesters of general chemistry, organic chemistry, calculus, and physics.
%
% --- MATERIALS ---
%
\section{Materials}
The following items are required for this course. Additional details can be found on the course website under the Materials section.
\begin{enumerate}
\item \emph{Physical Chemistry} by Thomas Engel and Phillip Reid, Third Edition, Prentice Hall (2013).  ISBN-13: 978-0321812001 \marginnote{-0.3in}{\href{http://goo.gl/qv2ndf}{http://goo.gl/qv2ndf}}
\item \emph{Quanta, Matter, and Change} by Peter Atkins, Julio de Paula, and Ronald Friedman, First Edition, W.H. Freeman and Company (2008). ISBN-13: 978-0716761174. \marginnote{-0.48in}{\href{http://goo.gl/0UzdHl}{http://goo.gl/0UzdHl}}
\item Wolfram's \emph{Mathematica}.
\item A working scientific calculator.
\end{enumerate}
%
% --- TIME REQUIREMENT ---
%
\section{Time Requirement}
Each of the two weekly course meetings is 110 minutes in length, requiring a total of 3.67 hours per week. Review meetings will, occasionally, add another hour per week. For each regular class meeting you can expect to spend at least 2 hours outside of class reading, studying, and working exercises. 
%
% --- CONTENT ---
%
\section{Content}
This course is an introduction to quantum mechanics, spectroscopy, and bonding applied to understanding chemical systems. The course topics are listed below. Topics are not necessarily covered in this order, and not all topics are covered with equal depth.
\begin{table}[h]
%\caption{\sffamily Topics covered in Chemistry 350.}
%\label{tab:topics}
\renewcommand{\arraystretch}{1}
\begin{tabular}{l|l} \toprule
Differential Equations & Atomic Spectroscopy \\
Beginnings of Quantum Mechanics & Computational Chemistry \\
Postulates of Quantum Mechanics & Molecular Spectroscopy \\
Schr\"{o}dinger's Wave Equation & Infrared Spectroscopy \\
Particle in a Box & Microwave Spectroscopy \\ 
Tunneling & Electronic Spectroscopy \\ 
Harmonic Oscillator & nmr Spectroscopy \\
Rigid Rotor & Variational Theory \\
Hydrogen Atom & Perturbation Theory \\
Multi-Electron Atoms &  Bonding in Diatomic Molecules \\
Atomic Term Symbols &  Bonding in Polyatomic Molecules \\
\bottomrule
\end{tabular}
\end{table}
%
% --- GOALS ---
%
\section{Goals}
Quantum \marginnote{-0.1in}{Quantum mechanics gives us the ability to predict chemistry without doing any wet laboratory work. What's not to like about that?} mechanics is the theoretical framework that organizes chemistry, starting with the periodic table. It allows us to understand the macroscopic behavior of matter in terms of the structure and interaction of individual atoms and molecules. It allows us to make sense of the interaction between matter and electromagnetic radiation. Its wide-ranging applicability makes it impossible to avoid quantum mechanics in any modern study of chemistry. It also happens to be a very cool achievement of the human mind. 

I want to share some of this with you. During this course I want to 
\begin{inparaenum}[\bfseries (a\upshape)]
\item provide you with the tools to understand how quantum mechanics applies to chemical systems;
\item provide you with the tools to understand quantum mechanics in the literature of any chemical subject;
\item show you the beauty and regularity of chemistry and Nature that emerge when you understand the most fundamental concepts that underlie everything you have ever experienced and everything you will ever experience;
\item provide you with the tools necessary to form a realistic mental image of the microscopic world of atoms and molecules;
\item further develop your ability to analyze and interpret experimental results in terms of the microscopic structure and dynamics of atoms and molecules; and
\item help you understand how we know what we know.
\end{inparaenum}

I also want you to move closer to being a professional in this course. Because of this, additional goals of this course are to
\begin{inparaenum}[\bfseries (a\upshape)]
\item improve your problem solving skills;
\item improve your ability to formulate descriptions of the physical and chemical world in terms of mathematical models; and
\item improve your ability to make clear and precise quantitative arguments.
\end{inparaenum}
%
% --- LEARNING OBJECTIVES ---
%
\section{Learning Objectives}
Learning objectives are things you should be able to do at the end of the course. For each topic covered in this course, you are provided with a list of learning objectives called Be Able Tos, or BATS for short. A complete list of BATS for each course topic can be found on the course website. These BATS are very detailed (granular), and it might be difficult to see the overall objectives of the course from them.

A higher level (less granular) list of learning objectives might be helpful. At the end of the course, you should be able to
\begin{inparaenum}[\bfseries (a\upshape)]
\item describe the difference between classical and quantum mechanics; 
\item describe the experiments that led to the development of quantum mechanics;
\item define wave-particle duality;
\item explain what quantization is and its origin;
\item solve the Schr\"odinger equation for a free particle and particle in a box;
\item use approximation techniques to solve the Schr\"odinger equation;
\item use molecular orbital theory to predict bonding in small molecules; and
\item predict and explain the outcome of electromagnetic radiation interacting with matter. 
\end{inparaenum}
%
% --- WEEKLY ROUTINE ---
%
\section{Weekly Routine}
Tuesday \marginnote{-0.1in}{The course website contains detailed schedules and lists of assignments.} and Thursday meetings are primarily dedicated to lecture. Ungraded exercises are assigned for each each topic, and you should work those on a regular basis. There are two exams, three micro exams, and one final exam.  
%
% --- THINGS I GRADE ---
%
\section{Things I Grade}
You and I determine your progress in this course using scores derived from evaluating the quality and accuracy of your answers to questions posed in micro exams, exams, and a final exam. This section provides details for each of these evaluations
%
% --- Micro Exams ---
%
\subsection{Micro Exams}
A micro exam is given at the start of a class meeting. Exact dates are in the course schedule. You are given $50$ minutes to complete a micro exam. The topics covered on each micro exam are listed in the course schedule. Solutions to micro exams are posted on the course website after they are graded. Each micro exam is worth 50 points.
%
% --- Exams ---
%
\subsection{Exams}
There are two exams during the semester. All exams are cumulative. Solutions for exams are posted on the course website after all exams are graded. See the course website for the exam schedule. Each exam is weighted equally at 150 points each. 
%
% --- Final Exam ---
%
\subsection{Final Exam}
The \marginnote{-0.1in}{Make plans to be on campus through the end of the final exam.} final exam is given on the day and at the time specified by the Registrar. There are no exceptions. The final exam is three hours in length, and it is cumulative over the entire semester. No solutions to the final exam are posted, but you may review the grading by making an appointment with me. You may not take possession of the exam.
%
% --- THINGS I DO NOT GRADE ---
%
\section{Things I Do Not Grade}
Each topic is accompanied by a set of practice exercises. These exercises are recommended, but I never collect and grade them. It is essential that you work (or at least review) as many of the practice exercises as possible. Solutions for practice exercises from the course textbooks are available outside my office. Solutions for other practice exercises can be found on the course website.
%
% --- COURSE SCORE ---
%
\section{Course Score}
Your \marginnote{-0.1in}{Weights of Evaluations \\ Micro Exams: \(35\%\) \\ Exams: \(35\%\) \\ Final Exam: \(30\%\)}course score is a weighted average of the scores you earn on micro exams and exams. These scores are weighted at \(35\%\) for micro exams, \(35\%\) for exams, and \(30\%\) for the final exam.  
%
% --- LETTER GRADE ---
%
\section{Letter Grade}
Letter grades are assigned at the end of the course according to the minimum course score requirements listed in Table~\ref{tab:lettergrades}. Course scores below $55\%$ are considered failing. Please see the course website for more detail about how your course letter grade is determined. 
\begin{table}[h]
\caption{\sffamily Minimum course score necessary for each letter grade.}
\label{tab:lettergrades}
\begin{tabular}{cl||cl} \toprule
\textbf{Minimum Score} & \textbf{Letter Grade} & \textbf{Minimum Score} & \textbf{Letter Grade} \\ \hline
97 & \hspace{0.3in}A$+$ & 72 & \hspace{0.3in}C$+$ \\
88 & \hspace{0.3in}A & 68 & \hspace{0.3in}C \\
85 & \hspace{0.3in}A$-$ & 65 & \hspace{0.3in}C$-$ \\
82 & \hspace{0.3in}B$+$ & 62 & \hspace{0.3in}D$+$ \\
78 & \hspace{0.3in}B & 58 & \hspace{0.3in}D \\
75 & \hspace{0.3in}B$-$ & 55 & \hspace{0.3in}D$-$ \\
\bottomrule
\end{tabular}
\end{table}
%
% --- ADDITIONAL INFORMATION ---
%
\section{Additional Information}
Please see the course website for additional information such as the detailed the course schedule, office hours, granular course objectives, suggestions for success, my expectations of you, grade advice, and policies on assignment of letter grades. The website also contains detailed policies on attendance, late submission of work, email communication, disability accommodations, and classroom behavior. All material found on the course website is an extension of this syllabus and constitutes my contract with you for this course. 
%
% --- LICENSE AND SOURCE CODE ---
%
\section{License and Source Code}
\copyright\ 2016 by Dale J. Brugh. Course Syllabus for Chemistry 350 (2016) is made available under a Creative Commons Attribution-ShareAlike 4.0 International License (CC BY-SA 4.0). See \href{https://creativecommons.org/licenses/by-sa/4.0/}{https://creativecommons.org/licenses/by-sa/4.0/} for details about your rights under this license. The \LaTeX\ source code is available under the same license from Github at \href{https://github.com/djbrugh/chem350}{https://github.com/djbrugh/chem350}. \marginpar{\vspace{-0.85in}\hfill\noindent\href{https://creativecommons.org/licenses/by-sa/4.0/}{\includegraphics[width=1in]{figs/cc-by-sa.pdf}}}
%
% --- END DOCUMENT ---
%
\end{document}
